% 乘积法则（综述）
% license CCBYSA3
% type Wiki

本文根据 CC-BY-SA 协议转载翻译自维基百科\href{https://en.wikipedia.org/wiki/Product_rule}{相关文章}

在微积分中，乘积法则（或称莱布尼茨法则[2]、莱布尼茨乘积法则）是一个用于求两个或多个函数乘积的导数的公式。对于两个函数，可以用拉格朗日记号表示为：
$$
(u \cdot v)' = u' \cdot v + u \cdot v',~
$$
或用莱布尼茨记号表示为：
$$
\frac{d}{dx}(u \cdot v) = \frac{du}{dx} \cdot v + u \cdot \frac{dv}{dx}.~
$$
该法则可以推广到三个或更多函数的乘积，推广到乘积的高阶导数规则，以及其他语境中。
\subsection{发现}
这一法则的发现通常归功于 戈特弗里德·莱布尼茨，他使用“无穷小”（现代微分的前身）进行了论证。[3]（然而，莱布尼茨论文的译者 J. M. Child[4] 认为应归功于艾萨克·巴罗。）以下是莱布尼茨的推理：[5]设$u$和$v$是函数。则$d(uv)$等于两个相继$uv$之差；其中一个是$uv$，另一个是$(u+du)(v+dv)$。于是：
$$
d(u \cdot v) = (u+du)\cdot(v+dv) - u\cdot v 
= u \cdot dv + v \cdot du + du \cdot dv.~
$$
由于项$du \cdot dv$相较于$du$和$dv$是“可忽略的”，莱布尼茨得出结论：
$$
d(u \cdot v) = v \cdot du + u \cdot dv,~
$$
这确实就是乘积法则的微分形式。

如果两边同时除以微分$dx$，就得到：
$$
\frac{d}{dx}(u \cdot v) = v \cdot \frac{du}{dx} + u \cdot \frac{dv}{dx},~
$$
这同样可以用拉格朗日记号写作：
$$
(u \cdot v)' = v \cdot u' + u \cdot v'.~
$$
\subsubsection{最初的证明}
莱布尼茨和牛顿都给出了乘积法则的证明，但以现代标准来看并不严格。
莱布尼茨借助“无穷小量”进行推理，把积解释为矩形的面积；而牛顿则使用“流动量”的概念进行推理。[6][7]
\subsection{例子}
\begin{itemize}
\item 假设我们要求函数$f(x) = x^{2}\sin(x)$的导数。应用乘积法则，可以得到$f'(x) = 2x \cdot \sin(x) + x^{2}\cos(x)$，因为$x^{2}$的导数是 $2x$，而正弦函数的导数是余弦函数。
\item 乘积法则的一个特殊情形是常数倍法则：如果 $c$ 是一个常数，而$f(x)$是一个可微函数，那么$(cf)(x) = c \cdot f(x)$也是可微的，其导数为$(cf)'(x) = c \cdot f'(x)$。这一结果可以直接从乘积法则推出，因为任何常数的导数都是零。
\item 这一点结合求导的和法则，表明微分运算是线性的。此外，分部积分公式就是由乘积法则推导出来的，商法则（弱形式）也是如此。所谓“弱形式”，指的是它并没有证明商一定可微，而只是指出：如果可微，其导数是什么。
\end{itemize}
\subsection{证明}
\subsubsection{导数的极限定义}
设$h(x) = f(x)g(x)$，并且假设$f$和$g$在$x$处可导。我们要证明$h$在 $x$ 处也可导，且其导数为$h'(x) = f'(x)g(x) + f(x)g'(x)$.为此，在分子中加上并减去一项$f(x)g(x+\Delta x) - f(x)g(x+\Delta x)$（其值为零，因此不改变原式），从而可以因式分解，然后利用极限的性质来处理。
$$
\begin{aligned}
h'(x) &= \lim_{\Delta x \to 0} \frac{h(x+\Delta x) - h(x)}{\Delta x} \\[6pt]
&= \lim_{\Delta x \to 0} \frac{f(x+\Delta x)g(x+\Delta x) - f(x)g(x)}{\Delta x} \\[6pt]
&= \lim_{\Delta x \to 0} \frac{f(x+\Delta x)g(x+\Delta x) - f(x)g(x+\Delta x) + f(x)g(x+\Delta x) - f(x)g(x)}{\Delta x} \\[6pt]
&= \lim_{\Delta x \to 0} \frac{\big(f(x+\Delta x)-f(x)\big)\cdot g(x+\Delta x) + f(x)\cdot \big(g(x+\Delta x)-g(x)\big)}{\Delta x} \\[6pt]
&= \lim_{\Delta x \to 0} \frac{f(x+\Delta x)-f(x)}{\Delta x} \cdot \lim_{\Delta x \to 0} g(x+\Delta x) \\ 
&\quad + \lim_{\Delta x \to 0} f(x) \cdot \lim_{\Delta x \to 0} \frac{g(x+\Delta x)-g(x)}{\Delta x} \\[6pt]
&= f'(x)g(x) + f(x)g'(x).
\end{aligned}~
$$
其中$\lim_{\Delta x \to 0} g(x+\Delta x) = g(x)$是因为可导函数必定连续。
\subsubsection{线性近似}
根据定义，如果$f,g:\mathbb{R}\to \mathbb{R}$在$x$处可导，则我们可以写成线性近似式：
$$
f(x+h) = f(x) + f'(x)h + \varepsilon_{1}(h),~
$$
$$
g(x+h) = g(x) + g'(x)h + \varepsilon_{2}(h),~
$$
其中误差项相对于$h$是高阶无穷小：
$$
\lim_{h\to 0}\frac{\varepsilon_{1}(h)}{h} = \lim_{h\to 0}\frac{\varepsilon_{2}(h)}{h} = 0,~
$$
也记作 $\varepsilon_{1}, \varepsilon_{2} \sim o(h)$。

于是：
$$
\begin{aligned}
f(x+h)g(x+h) - f(x)g(x) 
&= \big(f(x) + f'(x)h + \varepsilon_{1}(h)\big)\big(g(x) + g'(x)h + \varepsilon_{2}(h)\big) - f(x)g(x) \\[6pt]
&= f(x)g(x) + f'(x)g(x)h + f(x)g'(x)h - f(x)g(x) + \text{误差项} \\[6pt]
&= f'(x)g(x)h + f(x)g'(x)h + o(h).
\end{aligned}~
$$
这里的“误差项”包括诸如
$$
f(x)\varepsilon_{2}(h), \quad f'(x)g'(x)h^{2}, \quad h f'(x)\varepsilon_{1}(h),~
$$
这些量显然都是$o(h)$级别的。两边同时除以$h$，再令$h \to 0$，即可得到结果：$(fg)'(x) = f'(x)g(x) + f(x)g'(x)$.
\subsubsection{四分平方法}
这个证明利用了链式法则和四分平方函数$q(x) = \tfrac{1}{4}x^{2}, \quad q'(x) = \tfrac{1}{2}x$我们有：
$$
uv = q(u+v) - q(u-v),~
$$
对两边同时求导：
$$
\begin{aligned}
f' &= q'(u+v)(u' + v') - q'(u-v)(u' - v') \\[6pt]
&= \tfrac{1}{2}(u+v)(u' + v') - \tfrac{1}{2}(u-v)(u' - v') \\[6pt]
&= \tfrac{1}{2}(uu' + vu' + uv' + vv') - \tfrac{1}{2}(uu' - vu' - uv' + vv') \\[6pt]
&= vu' + uv'.
\end{aligned}~
$$
\subsubsection{多变量链式法则}
乘积法则可以看作是多元链式法则在乘法函数$m(u, v) = uv$上的一种特殊情况：
$$
\frac{d(uv)}{dx} = \frac{\partial (uv)}{\partial u}\frac{du}{dx} + \frac{\partial (uv)}{\partial v}\frac{dv}{dx} = v\frac{du}{dx} + u\frac{dv}{dx}.~
$$
\subsubsection{非标准分析}
设$u$和$v$是关于$x$的连续函数，并且$dx, du, dv$是非标准分析框架下（具体来说是超实数体系中的）无穷小。记$\mathrm{st}$为标准部函数，它将有限超实数映射为与之无限接近的实数。于是有：
$$
\begin{aligned}
\frac{d(uv)}{dx}
&= \operatorname{st}\!\left(\frac{(u+du)(v+dv)-uv}{dx}\right)\\
&= \operatorname{st}\!\left(\frac{uv+u\cdot dv+v\cdot du+du\cdot dv - uv}{dx}\right)\\
&= \operatorname{st}\!\left(\frac{u\cdot dv+v\cdot du+du\cdot dv}{dx}\right)\\
&= \operatorname{st}\!\left(u\frac{dv}{dx}+(v+dv)\frac{du}{dx}\right)\\
&= u\frac{dv}{dx}+v\frac{du}{dx}.
\end{aligned}~
$$
这实际上就是莱布尼茨当时的证明方式，本质上依赖于同类项超越法则，只不过这里用的是“标准部”函数来代替。
\subsubsection{光滑无穷小分析}
在Lawvere的无穷小方法框架下，设$dx$是一个平方为零的无穷小量。于是有$du = u'\,dx, \quad dv = v'\,dx$,因此：
$$
\begin{aligned}
d(uv) &= (u+du)(v+dv) - uv \\
&= uv + u\cdot dv + v\cdot du + du\cdot dv - uv \\
&= u\cdot dv + v\cdot du + du\cdot dv \\
&= u\cdot dv + v\cdot du
\end{aligned}~
$$
因为$du\,dv = u'v'(dx)^2 = 0$.再除以$dx$，得到：$\frac{d(uv)}{dx} = u\frac{dv}{dx} + v\frac{du}{dx}$,或者等价地：$(uv)' = u\cdot v' + v\cdot u'$.
\subsubsection{对数微分法}
设$h(x) = f(x)g(x)$.对等式两边取绝对值并取自然对数：
$$
\ln |h(x)| = \ln |f(x)g(x)|.~
$$
利用绝对值和对数的性质：
$$
\ln |h(x)| = \ln |f(x)| + \ln |g(x)|.~
$$
接着对等式两边取对数导数，然后解出 $h'(x)$：
$$
\frac{h'(x)}{h(x)} = \frac{f'(x)}{f(x)} + \frac{g'(x)}{g(x)}.~
$$
再解出$h'(x)$，并将$h(x)=f(x)g(x)$ 代回：
$$
\begin{aligned}
h'(x) &= h(x)\left(\frac{f'(x)}{f(x)} + \frac{g'(x)}{g(x)}\right) \\
&= f(x)g(x)\left(\frac{f'(x)}{f(x)} + \frac{g'(x)}{g(x)}\right) \\
&= f'(x)g(x) + f(x)g'(x).
\end{aligned}~
$$
\textbf{注意：}在对数微分时必须对函数取绝对值，因为对数函数 $\ln u$ 仅在 $u > 0$ 时为实值。但这仍然成立，因为$\frac{d}{dx}(\ln |u|) = \frac{u'}{u}$,这就合理化了对函数取绝对值的步骤。
\subsection{推广}
\subsubsection{多个因子的乘积}
乘积法则可以推广到两个以上因子的情形。例如，对三个因子有：
$$
\frac{d(uvw)}{dx} = \frac{du}{dx}vw + u\frac{dv}{dx}w + uv\frac{dw}{dx}.~
$$
对于函数族 $f_{1}, \dots , f_{k}$，有：
$$
\frac{d}{dx}\left[\prod_{i=1}^{k} f_{i}(x)\right] 
= \sum_{i=1}^{k}\left(\left(\frac{d}{dx}f_{i}(x)\right)\prod_{j=1, j\neq i}^{k} f_{j}(x)\right) 
= \left(\prod_{i=1}^{k} f_{i}(x)\right)\left(\sum_{i=1}^{k}\frac{f'_{i}(x)}{f_{i}(x)}\right).~
$$
对数导数能给出更简洁的形式，并且提供一种不涉及递归的直接证明。函数$f$的对数导数（记作$\operatorname{Logder}(f)$）定义为该函数对数的导数：
$$
\operatorname{Logder}(f) = \frac{f'}{f}.~
$$
利用“乘积的对数等于因子对数的和”这一性质，结合导数的和法则，可以直接得到：
$$
\operatorname{Logder}(f_{1}\cdots f_{k}) = \sum_{i=1}^{k}\operatorname{Logder}(f_{i}).~
$$
将等式两边同时乘以 $\prod_{i=1}^{k} f_{i}$，就得到了上面关于乘积求导的表达式。
\subsubsection{高阶导数}
乘积法则也可以推广为两个因子乘积的 n 阶导数的一般莱布尼茨公式，其形式可由二项式定理展开得到：
$$
d^{n}(uv) = \sum_{k=0}^{n} {n \choose k} \cdot d^{(n-k)}(u) \cdot d^{(k)}(v).~
$$
应用于某一点 $x$，上式可写为：
$$
(uv)^{(n)}(x) = \sum_{k=0}^{n} {n \choose k} \cdot u^{(n-k)}(x) \cdot v^{(k)}(x).~
$$
此外，对于任意多个因子的乘积的 n 阶导数，有类似的公式，不过需要用到**多项式系数：
$$
\left(\prod_{i=1}^{k} f_{i}\right)^{(n)} 
= \sum_{j_{1}+j_{2}+\cdots+j_{k}=n} 
{n \choose j_{1}, j_{2}, \ldots, j_{k}} 
\prod_{i=1}^{k} f_{i}^{(j_{i})}.~
$$
这里，${n \choose j_{1}, j_{2}, \ldots, j_{k}}$ 表示多项式系数。 
\subsubsection{高阶偏导数}
对于偏导数，我们有【8】：
$$
\frac{\partial^{n}}{\partial x_{1}\cdots \partial x_{n}}(uv) 
= \sum_{S} 
\frac{\partial^{|S|}u}{\prod_{i\in S}\partial x_{i}}
\cdot
\frac{\partial^{\,n-|S|}v}{\prod_{i\notin S}\partial x_{i}}~
$$
其中，指标$S$遍历集合$\{1,\ldots,n\}$的所有$2^n$个子集，$|S|$表示子集的基数。

例如，当 $n=3$ 时：
$$
\begin{aligned}
\frac{\partial^{3}}{\partial x_{1}\partial x_{2}\partial x_{3}}(uv) 
={}& u \cdot \frac{\partial^{3}v}{\partial x_{1}\partial x_{2}\partial x_{3}}
+ \frac{\partial u}{\partial x_{1}} \cdot \frac{\partial^{2}v}{\partial x_{2}\partial x_{3}}
+ \frac{\partial u}{\partial x_{2}} \cdot \frac{\partial^{2}v}{\partial x_{1}\partial x_{3}}
+ \frac{\partial u}{\partial x_{3}} \cdot \frac{\partial^{2}v}{\partial x_{1}\partial x_{2}} \\
&+ \frac{\partial^{2}u}{\partial x_{1}\partial x_{2}} \cdot \frac{\partial v}{\partial x_{3}}
+ \frac{\partial^{2}u}{\partial x_{1}\partial x_{3}} \cdot \frac{\partial v}{\partial x_{2}}
+ \frac{\partial^{2}u}{\partial x_{2}\partial x_{3}} \cdot \frac{\partial v}{\partial x_{1}}
+ \frac{\partial^{3}u}{\partial x_{1}\partial x_{2}\partial x_{3}} \cdot v.
\end{aligned}~
$$
\subsubsection{巴拿赫空间}
假设$X, Y, Z$是巴拿赫空间（其中也包括欧几里得空间），并且$B : X \times Y \to Z$是一个连续双线性算子。那么$B$是可微的，并且它在点$(x, y)\in X \times Y$的导数是一个线性映射$D_{(x,y)}B : X \times Y \to Z$，其形式为：
$$
\bigl(D_{(x,y)}B\bigr)(u,v) = B(u,y) + B(x,v) \quad \forall (u,v)\in X\times Y.~
$$
这一结果可以推广【9】到更一般的拓扑向量空间。
\subsubsection{在向量分析中}
乘积法则可以扩展到 $\mathbb{R}^n$ 上向量函数的各种乘积运算：[10]
\begin{itemize}
\item 数量乘法：
$$
(f \cdot g)' = f' \cdot g + f \cdot g'~
$$
\item 点积：
$$
(\mathbf{f} \cdot \mathbf{g})' = \mathbf{f}' \cdot \mathbf{g} + \mathbf{f} \cdot \mathbf{g}'~
$$
\item 叉积（cross product，在 $\mathbb{R}^3$ 上）：
$$
(\mathbf{f} \times \mathbf{g})' = \mathbf{f}' \times \mathbf{g} + \mathbf{f} \times \mathbf{g}'~
$$
\end{itemize}
若 $f$ 和 $g$ 是标量场，那么在梯度（gradient）意义下也存在一个乘积法则：
$$
\nabla(f \cdot g) = (\nabla f)\, g + f \, (\nabla g)~
$$
这样的法则对任何连续双线性乘积运算都成立。设$B : X \times Y \to Z$是向量空间之间的一个连续双线性映射，且$f$和$g$分别是到$X$和$Y$的可微函数。利用导数的极限定义时，乘法所使用的唯一性质就是它的连续性和双线性。因此，对于任何连续双线性运算，都有：
$$
H(f,g)' = H(f', g) + H(f, g').~
$$
这也是 Banach 空间中关于双线性映射乘积法则的一个特例。
\subsubsection{抽象代数与微分几何中的导子}
在抽象代数中，乘积法则是“导子”（derivation）的定义性质。在这种术语下，乘积法则表明导数算子是函数上的一个导子。

在微分几何中，流形$M$上一点$p$的切向量可以抽象地定义为一个作用在实值函数上的算子，其行为类似于在$p$点的方向导数。换句话说，切向量是一个线性泛函 $v$，并且它满足导子性质：
$$
v(fg) = v(f)\, g(p) + f(p)\, v(g).~
$$
进一步推广（并对偶化）向量分析的公式到 $n$ 维流形 $M$，我们可以取次数为 $k$ 和 $\ell$ 的微分形式 $\alpha \in \Omega^k(M)$，$\beta \in \Omega^\ell(M)$，它们之间有外积运算：
$$
\alpha \wedge \beta \in \Omega^{k+\ell}(M),~
$$
以及外微分算子：
$$
d : \Omega^m(M) \to \Omega^{m+1}(M).~
$$
此时，成立分次 Leibniz 法则：
$$
d(\alpha \wedge \beta) = d\alpha \wedge \beta + (-1)^k \alpha \wedge d\beta.~
$$
\subsection{应用}
乘积法则的一个应用是证明：
$$
\frac{d}{dx} x^n = n x^{n-1}~
$$
当 $n$ 是正整数时成立。（事实上，该法则在 $n$ 不是正数甚至不是整数时依然成立，但那时的证明必须依赖其他方法。）证明方法是对指数 $n$ 使用数学归纳法。基例：若 $n=0$，则 $x^n$ 是常数，而 $n x^{n-1} = 0$。在此情况下，法则成立，因为常函数的导数为 0。归纳步骤：假设该法则对某个特定的指数 $n$ 成立，则对下一个指数 $n+1$，有：
$$
\begin{aligned}
\frac{d}{dx} x^{n+1} 
&= \frac{d}{dx}(x^n \cdot x) \\
&= x \frac{d}{dx} x^n + x^n \frac{d}{dx} x \quad (\text{此处用到了乘积法则})\\
&= x \left( n x^{n-1} \right) + x^n \cdot 1 \quad (\text{此处用了归纳假设})\\
&= (n+1) x^n.~
\end{aligned}~
$$
因此，如果该命题对 $n$ 成立，那么它对 $n+1$ 也成立，从而对所有自然数 $n$ 都成立。
\subsection{参见}
\begin{itemize}
\item 积分的求导 —— 数学中的一个问题
\item 三角函数的求导 —— 求三角函数导数的数学过程
\item 求导法则 —— 计算函数导数的规则
\item 分布（数学） —— 推广函数概念的数学术语
\item 广义莱布尼茨法则 —— 微积分中乘积法则的推广
\item 分部积分法 —— 微积分中的一种数学方法
\item 反函数与求导 —— 反函数导数的公式
\item 求导的线性性 —— 微积分的一个性质
\item 幂函数求导法则 —— 单项式多项式的求导方法
\item 商法则 —— 函数比值的导数公式
\item 导数表 —— 计算函数导数的规则汇编
\item 向量分析恒等式 —— 数学恒等式
\end{itemize}
\subsection{参考文献}

* 注：这是自 17 世纪以来常见的一幅图，与 James Stewart《Calculus: Early Transcendentals》第 7 版第 185 页“乘积法则的几何意义”一节中的插图基本相同。
* “Leibniz rule – 《数学百科全书》”。
* Michelle Cirillo（2007 年 8 月），“Humanizing Calculus”，*The Mathematics Teacher*，第 101 卷（第 1 期）：23–27。doi:10.5951/MT.101.1.0023。
* Leibniz, G. W. (2005) \[1920], *The Early Mathematical Manuscripts of Leibniz* (PDF)，J.M. Child 翻译，Dover 出版，第 28 页，脚注 58，ISBN 978-0-486-44596-0。
* Leibniz, G. W. (2005) \[1920], *The Early Mathematical Manuscripts of Leibniz* (PDF)，J.M. Child 翻译，Dover 出版，第 143 页，ISBN 978-0-486-44596-0。
* Eugene Boman, and Robert Rogers. *Real Analysis* \[在线教材] [https://math.libretexts.org/Bookshelves/Analysis/Real\_Analysis\_(Boman\_and\_Rogers)/02\%3A\_Calculus\_in\_the\_17th\_and\_18th\_Centuries/2.01\%3A\_Newton\_and\_Leibniz\_Get\_Started](https://math.libretexts.org/Bookshelves/Analysis/Real_Analysis_\%28Boman_and_Rogers\%29/02\%3A_Calculus_in_the_17th_and_18th_Centuries/2.01\%3A_Newton_and_Leibniz_Get_Started)
* “A Story of Real Analysis” (PDF). 原始 PDF 于 2025 年 3 月 11 日存档。
* Micheal Hardy（2006 年 1 月），“Combinatorics of Partial Derivatives” (PDF)，*The Electronic Journal of Combinatorics*，第 13 卷。arXiv\:math/0601149。Bibcode:2006math......1149H。
* Kriegl, Andreas; Michor, Peter (1997). *The Convenient Setting of Global Analysis* (PDF). American Mathematical Society，第 59 页。ISBN 0-8218-0780-3。
* Stewart, James (2016), *Calculus*（第 8 版），Cengage，第 13.2 节。

---

要不要我也帮你整理成 **中英文对照参考文献格式**（比如像学术论文那样排版），方便引用？
